% ==============================================================================
% Fichier : chapitres/04_ingenierie_architecture.tex
% Livre IV : Ingénierie et Architecture de Sécurité
% ==============================================================================

\chapter{Ingénierie et Architecture de Sécurité}
\label{chap:ingenierie_archi}

\begin{boxobjectifs}
\textbf{Objectifs du Livre~IV :} Donner les éléments liés aux modèles de sécurité, aux fondements de la cryptographie (symétrique, asymétrique, hachage), à la sécurité physique et à la virtualisation. À l'issue, l'étudiant sera capable d'implémenter les chiffrements classiques et de concevoir une architecture de sécurité défense en profondeur.
\end{boxobjectifs}

% ------------------------------------------------------------------------------
\section{Modèles de Sécurité}
% ------------------------------------------------------------------------------

\subsection{Modèle de Bell-LaPadula (Confidentialité)}

\begin{boxdef}
Modèle orienté \textbf{confidentialité}, utilisé dans les systèmes militaires. Repose sur deux règles :
\begin{itemize}
    \item \textbf{No Read Up (NRU) :} Un sujet ne peut pas lire une ressource de niveau de classification supérieur au sien.
    \item \textbf{No Write Down (NWD) :} Un sujet ne peut pas écrire dans une ressource de niveau inférieur au sien (pour éviter la fuite d'informations classifiées).
\end{itemize}
\end{boxdef}

\subsection{Modèle de Biba (Intégrité)}
\begin{boxdef}
Modèle orienté \textbf{intégrité}, inverse de Bell-LaPadula :
\begin{itemize}
    \item \textbf{No Write Up :} Un sujet ne peut pas écrire dans une ressource de niveau d'intégrité supérieur.
    \item \textbf{No Read Down :} Un sujet ne peut pas lire une ressource de niveau d'intégrité inférieur.
\end{itemize}
\end{boxdef}

\subsection{Chinese Wall (Brewer-Nash)}
Modèle dynamique évitant les \textbf{conflits d'intérêts}. Un consultant ayant accédé aux données d'une entreprise A ne peut plus accéder aux données d'une entreprise concurrente B.

% ------------------------------------------------------------------------------
\section{Fondamentaux de la Conception d'Architecture de Sécurité}
% ------------------------------------------------------------------------------

\subsection{Principes de Base}

\begin{itemize}
    \item \textbf{Défense en profondeur :} Multiplier les couches de sécurité indépendantes. Si une couche est compromise, les suivantes assurent la protection.
    \item \textbf{Moindre privilège :} N'accorder que les droits strictement nécessaires à l'accomplissement d'une tâche.
    \item \textbf{Séparation des fonctions :} Aucun individu ne doit avoir le contrôle complet sur un processus critique (principe des «~quatre yeux~»).
    \item \textbf{Sécurité par défaut :} Les systèmes doivent être sécurisés dans leur configuration initiale.
    \item \textbf{Fail Secure :} En cas de panne, le système doit basculer dans un état sécurisé (accès refusé plutôt qu'accordé).
\end{itemize}

\subsection{Segmentation Réseau et Défense en Profondeur}

\begin{center}
\begin{tikzpicture}[scale=0.85]
  \draw[thick, CouleurPrimaire, fill=CouleurPrimaire!5, rounded corners] (-5,-3) rectangle (5,3);
  \node[CouleurPrimaire, font=\small\bfseries] at (0,2.7) {Zone Externe (Internet)};

  \draw[thick, CouleurAccent, fill=CouleurAccent!5, rounded corners] (-3.5,-2) rectangle (3.5,2);
  \node[CouleurAccent, font=\small\bfseries] at (0,1.7) {DMZ (Serveurs Web, Mail)};

  \draw[thick, NavyBlue, fill=NavyBlue!5, rounded corners] (-2,-1) rectangle (2,1);
  \node[NavyBlue, font=\small\bfseries] at (0,0.7) {LAN Interne};
  \node[font=\tiny] at (0,0.1) {Serveurs applicatifs, BDD};
  \node[font=\tiny] at (0,-0.3) {Postes de travail};

  % Pares-feux
  \node[font=\tiny\bfseries, CouleurPrimaire] at (4,-2.5) {Pare-feu 1};
  \node[font=\tiny\bfseries, CouleurAccent]  at (3,-1.5) {Pare-feu 2};
\end{tikzpicture}
\captionof{figure}{Architecture Défense en Profondeur avec DMZ}
\end{center}

% ------------------------------------------------------------------------------
\section{Architecture Matérielle Sécurisée}
% ------------------------------------------------------------------------------

\subsection{Vulnérabilités et Menaces Matérielles}
\begin{itemize}
    \item Défauts de conception ou de fabrication (ex. : Spectre, Meltdown).
    \item Configurations incorrectes (BIOS/UEFI non protégés, ports USB ouverts).
    \item Menaces physiques : vol, sabotage, espionnage industriel.
\end{itemize}

\subsection{Mesures de Protection Matérielle}
\begin{itemize}
    \item \textbf{TPM (Trusted Platform Module) :} Puce sécurisée pour le stockage de clés de chiffrement et la vérification de l'intégrité du démarrage.
    \item \textbf{Contrôle d'accès physique :} Salles serveurs verrouillées, racks sécurisés.
    \item \textbf{Surveillance de l'intégrité :} Journaux événements, IDS matériels.
    \item \textbf{Mises à jour de micrologiciel :} Application régulière des patchs BIOS/UEFI.
\end{itemize}

% ------------------------------------------------------------------------------
\section{Systèmes d'Exploitation et Architecture Logicielle}
% ------------------------------------------------------------------------------

\subsection{Vulnérabilités des Systèmes d'Exploitation}
\begin{itemize}
    \item Erreurs de programmation (buffer overflow, use-after-free).
    \item Droits d'accès inappropriés (comptes sur-privilégiés).
    \item Services non utilisés exposant des surfaces d'attaque inutiles.
    \item Mises à jour non appliquées (exploitation des CVE connus).
\end{itemize}

\subsection{Principes d'Architecture Logicielle Sécurisée}
\begin{itemize}
    \item \textbf{Modularité :} Composants indépendants et substituables.
    \item \textbf{Séparation des préoccupations :} Fonctions de sécurité isolées du code métier.
    \item \textbf{Abstraction :} Masquage des détails d'implémentation.
    \item \textbf{Évolutivité :} Capacité à gérer une charge croissante sans dégrader la sécurité.
\end{itemize}

% ------------------------------------------------------------------------------
\section{Virtualisation}
% ------------------------------------------------------------------------------

La virtualisation utilise un \termecle{hyperviseur} pour créer et gérer des machines virtuelles (VM). Elle introduit de nouvelles vulnérabilités :
\begin{itemize}
    \item \textbf{Attaques contre l'hyperviseur (VM Escape) :} Compromission de toutes les VM hébergées.
    \item \textbf{Attaques inter-VM (VM Hopping) :} Une VM compromise attaque les voisines.
    \item \textbf{Fuites de données mémoire} entre VM partageant le même hôte physique.
\end{itemize}

\textbf{Contre-mesures :} Sécurisation de l'hyperviseur, isolation des réseaux virtuels (VLAN), application stricte des politiques de groupes de sécurité.

% ------------------------------------------------------------------------------
\section{Introduction à la Cryptographie}
% ------------------------------------------------------------------------------

\subsection{Algorithmes Symétriques}

La même clé est utilisée pour chiffrer et déchiffrer. Alice et Bob doivent partager la clé secrète à l'avance.

\subsubsection{Chiffrement de César}
Substitution mono-alphabétique par décalage d'un entier $k$ ($0 \leq k \leq 25$) dans l'alphabet.

\[
C_i = (P_i + k) \mod 26 \qquad \text{(Chiffrement)}
\]
\[
P_i = (C_i - k) \mod 26 \qquad \text{(Déchiffrement)}
\]

\begin{boxlizbiz}
\textbf{Exemple :} $k=3$, message \texttt{BONJOUR} $\rightarrow$ \texttt{ERQMRXU}\\
\textit{Faiblesse :} 25 clés seulement, cassé par analyse de fréquence.
\end{boxlizbiz}

\subsubsection{Chiffrement de Vigenère}
Substitution poly-alphabétique utilisant une clé de longueur $n$. La lettre en position $i$ est décalée de la valeur numérique de la lettre de clé correspondante.

\[
C_i = (P_i + K_{i \mod n}) \mod 26
\]

\begin{boxlizbiz}
\textbf{Exemple :} Message \texttt{BONJOUR}, Clé \texttt{CLE} :
\begin{verbatim}
B O N J O U R
C L E C L E C
D Z R L Z Y T
\end{verbatim}
\textit{Faiblesse :} Cassé par le test de Kasiski si la clé est courte et répétée.
\end{boxlizbiz}

\subsubsection{Chiffrement de Hill}
Chiffrement par blocs utilisant l'algèbre matricielle sur $\mathbb{Z}_{26}$ :
\[
\mathbf{C} = \mathbf{K} \cdot \mathbf{P} \mod 26
\]
où $\mathbf{K}$ est la matrice-clé (inversible sur $\mathbb{Z}_{26}$) et $\mathbf{P}$ le vecteur plaintext.

\subsubsection{Algorithmes Modernes Symétriques}
\begin{itemize}
    \item \textbf{DES} (Data Encryption Standard, 1977) : 56 bits, obsolète.
    \item \textbf{3DES} : Triple DES, plus sûr mais lent.
    \item \textbf{AES} (Advanced Encryption Standard, 2001) : Standard actuel — 128, 192 ou 256 bits.
\end{itemize}

\subsection{Algorithmes Asymétriques}
Deux clés différentes : clé publique (chiffrement, vérification de signature) et clé privée (déchiffrement, signature). La clé publique peut être librement partagée.

\begin{itemize}
    \item \textbf{RSA :} Basé sur la difficulté de factorisation de grands entiers. Clés de 2048 ou 4096 bits.
    \item \textbf{ECDSA :} Basé sur les courbes elliptiques. Efficace pour les signatures.
    \item \textbf{Diffie-Hellman :} Protocole d'échange de clés sur un canal non sécurisé.
\end{itemize}

\begin{boxinfo}
\textbf{Utilisation hybride :} En pratique, on utilise la cryptographie asymétrique pour échanger une clé symétrique de session, puis AES pour chiffrer les données (TLS/SSL).
\end{boxinfo}

\subsection{Fonctions de Hachage}
Transforme des données de taille quelconque en empreinte de taille fixe. Propriétés : déterminisme, résistance aux préimages, résistance aux collisions, effet avalanche.

\begin{table}[H]
\centering
\begin{tabular}{llll}
    \toprule
    \textbf{Algorithme} & \textbf{Taille du hash} & \textbf{Statut} \\
    \midrule
    MD5    & 128 bits & Obsolète (collisions connues) \\
    SHA-1  & 160 bits & Déprécié \\
    SHA-256 & 256 bits & Recommandé \\
    SHA-3  & 256/512 bits & Standard actuel \\
    \bottomrule
\end{tabular}
\caption{Fonctions de hachage courantes}
\end{table}

\subsection{Cryptanalyse}
Science visant à analyser et casser les systèmes cryptographiques. Principales attaques :
\begin{itemize}
    \item \textbf{Force brute :} Essai exhaustif de toutes les clés possibles.
    \item \textbf{Analyse de fréquence :} Exploitation des redondances statistiques de la langue.
    \item \textbf{Attaque à texte clair connu :} Exploitation de paires (texte clair, texte chiffré) connues.
    \item \textbf{Attaque par collision :} Recherche de deux entrées distinctes produisant le même hash.
\end{itemize}

% ------------------------------------------------------------------------------
\section{Sécurité Physique}
% ------------------------------------------------------------------------------
\begin{itemize}
    \item Contrôle d'accès aux salles serveurs (badge, biométrie, sas).
    \item Surveillance vidéo et enregistrement des entrées.
    \item Protection contre les incendies (détection, suppression gaz FM200).
    \item Alimentation sans coupure (UPS) et générateurs.
    \item Protection contre les inondations, la chaleur et les interférences électromagnétiques.
\end{itemize}

% ------------------------------------------------------------------------------
\section{Évaluation des Acquis --- Livre~IV}
% ------------------------------------------------------------------------------

\begin{boxexo}
\textbf{Exercices :}
\begin{enumerate}
    \item Chiffrez le message \texttt{ENSPY} avec le chiffrement de César, $k=5$.
    \item Avec la clé \texttt{SEC}, chiffrez \texttt{AUDIT} avec Vigenère.
    \item Quelle est la différence fondamentale entre un algorithme symétrique et un algorithme asymétrique ?
    \item Expliquez le principe de défense en profondeur avec un schéma réseau commenté.
    \item Calculez le hash SHA-256 de la chaîne \texttt{Securite2024} (à faire en TP avec Python : \texttt{hashlib.sha256}).
\end{enumerate}

\textbf{Travail Pratique :}\\
Concevoir et développer un site web sécurisé implémentant le modèle Chinese Wall pour une entreprise de traitement de données. Le site doit : (1) authentifier les utilisateurs avec un hash sécurisé (bcrypt), (2) vérifier l'intégrité des données (SHA-256), (3) appliquer les restrictions Chinese Wall selon les projets affectés.
\end{boxexo}
